%------------------------------------------------------------------------------%
\documentclass[a4paper,12pt,fleqn]{article}

\usepackage{calculo_varias_variaveis-1}

\ShowAnswers

%------------------------------------------------------------------------------%
\begin{document}

\makeHeader{CFVVI}{Prova 2 -- Turma 7}{04/06/2025}

% 01
%------------------------------------------------------------------------------%
\exercise{25}
Considere que a equação
\(
  \ln(z) + xyz = y^2
\)
defina $z$ como função de $x$ e $y$, encontre \(\D{z}{x}\)
\clearpagequestiononly

\begin{answer}
  \begin{align*}
    \ln(z) + xyz                             & = y^2                   \\
    \D{}{x} \left( \ln(z) + xyz \right)      & = \D{}{x}y^2            \\
    \frac{1}{z} \D{z}{x} + yz + xy \D{z}{x}  & = 0                     \\
    \left( \frac{1}{z} + xy \right) \D{z}{x} & = -yz                   \\
    \frac{1+ xyz}{z} \; \D{z}{x}             & = -yz                   \\
    \D{z}{x}                                 & = \frac{-yz^2}{1 + xyz}
  \end{align*}
\end{answer}

% 02
%------------------------------------------------------------------------------%
\exercise{25}
Seja \( f(x, y) = \sqrt{4x + y^2} \),
use a aproximação linear de $f$ no ponto $(1, 2)$
para estimar $f(1,02; 1,98)$
\clearpagequestiononly

\begin{answer}
  Calculando as derivadas
  \begin{align*}
    \D{f}{x} & = \frac{2}{\sqrt{4x + y^2}} \\
    \D{f}{y} & = \frac{y}{\sqrt{4x + y^2}}
  \end{align*}
  Avaliando a função e as derivadas no ponto $(1, 2)$
  \begin{align*}
    f(1, 2)        & = \sqrt{4 + 4} = \sqrt{8} = 2\sqrt{2}      \\
    \D{f}{x}(1, 2) & = \frac{2}{2\sqrt{2}} = \frac{1}{\sqrt{2}} \\
    \D{f}{y}(1, 2) & = \frac{2}{2\sqrt{2}} = \frac{1}{\sqrt{2}}
  \end{align*}
  A aproximação linear é
  \begin{align*}
    L(x, y) & = f(x_0, y_0) + \D{f}{x}(x_0, y_0)(x-x_0) + \D{f}{y}(y-y_0) \\
            & = f(1, 2) + \D{f}{x}(1, 2)(x - 1) + \D{f}{y}(1, 2)(y - 2)   \\
            & = 2\sqrt{2} + \frac{1}{\sqrt{2}}(x - 1) + \frac{1}{\sqrt{2}}(y - 2) \\
            & = \frac{4}{\sqrt{2}} + \frac{x + y - 3}{\sqrt{2}} \\
            & = \frac{x + y + 1}{\sqrt{2}}
  \end{align*}
  \[
    L(1,02; 1,98) = \frac{1,02 + 1,98 + 1}{\sqrt{2}}
                  = \frac{4}{\sqrt{2}}
                  = 2\sqrt{2}
                  \approx 2.8284
  \]
  Comparando com o valor exato
  \[
    f(1,02; 1,98) = 2.8285
  \]
\end{answer}

% 03
%------------------------------------------------------------------------------%
\exercise{25}
Seja \( f(x, y) = \ln(x^2 + y^2 + 1) \)
\begin{enumerate}[label=\alph*)]
  \item Calcule o gradiente de $f$
  \item Calcule a derivada de $f$ no ponto $(1, 2)$ na direção do vetor $\vetor{-2}{1}$
\end{enumerate}
\clearpagequestiononly

\begin{answer}
  \begin{enumerate}[label=\alph*)]
    \item
    Vetor gradiente
    \[
    \nabla f(x, y)
    = \Vetor{\D{f}{x}}{\D{f}{y}}
    = \Vetor{\frac{2x}{x^2 + y^2 + 1}}{\frac{2y}{x^2 + y^2 + 1}}
    \]
    \item
    Gradiente no ponto $(1, 2)$
    \[
    \nabla f(x, y) (1, 2)
    = \Vetor{\frac{2}{6}}{\frac{4}{6}}
    = \Vetor{\frac{1}{3}}{\frac{2}{3}}
    \]
    Calculando um vetor unitário na direção de $v$
    \[
    u
    = \frac{v}{\norm{v}}
    = \frac{1}{\sqrt{(-2)^2 + 1^2}} \vetor{-2}{1}
    = \frac{1}{\sqrt{5}} \vetor{-2}{1}
    \]
    Derivada direcional
    \begin{align*}
      D_u f(1, 2)
      & = \nabla f(1, 2) \cdot u                       \\
      & = \Vetor{\frac{1}{3}}{\frac{2}{3}} \cdot \frac{1}{\sqrt{5}} \vetor{-2}{1} \\
      & = \frac{1}{\sqrt{5}} \left( \frac{-2}{3} + \frac{2}{3} \right) \\
      &  = 0
    \end{align*}
  \end{enumerate}
\end{answer}

% 04
%------------------------------------------------------------------------------%
\exercise{25}
Seja \( f(x, y, z) = x^2 \sin(yz) \),
calcule \( \DM{f}{y}{z} \) no ponto \( (1, \pi, 0) \).
Efetue as derivadas na ordem especificada pela notação.
\clearpagequestiononly

\begin{answer}
  \begin{align*}
    \D{f}{z}
    & = \D{}{z}\left(x^2 \sin(yz)\right) \\
    & = x^2 \D{}{z}\left(\sin(yz)\right) \\
    & = x^2 \cos(yz)\D{}{z}\left(yz\right) \\
    & = x^2 \cos(yz) y \\
    & = x^2 y \cos(yz)
  \end{align*}
  \begin{align*}
    \DM{f}{y}{z}
    & = \D{}{y}\left(x^2 y \cos(yz)\right) \\
    & = x^2 \D{}{y}\left(y \cos(yz)\right) \\
    & = x^2 \left[\D{y}{y}\cos(yz) + y\D{}{y}\cos(yz)\right] \\
    & = x^2 \left[\cos(yz) - y\sin(yz)\D{}{y}(yz)\right] \\
    & = x^2 \left[\cos(yz) - y\sin(yz)z\right] \\
    & = x^2 \left[\cos(yz) - yz\sin(yz)\right]
  \end{align*}
  \begin{align*}
    \DM{f}{y}{z}(1, \pi, 0)
    & = x^2 \left[\cos(yz) - yz\sin(yz)\right] \evalat{(1, \pi, 0)}{} \\
    & = 1^2 \left[\cos(\pi\times 0) - \pi\times 0\sin(\pi\times 0)\right] \\
    & = 1 \left[\cos(0) - 0\right] \\
    & = 1
  \end{align*}
\end{answer}

%------------------------------------------------------------------------------%
\end{document}
%------------------------------------------------------------------------------%

